\section{Optionaler Anhang: Matrix- und Robustheitstests}

Die Matrix ist ein kontrollierter Screening-Versuch, kein vollständiger
Vermessungsnachweis. Die historische STS-Auswertung umfasst den Batch
\texttt{matrix\_full\_pipe} und den Folge-Batch \texttt{matrix\_rest}. Die
sechs A-Läufe des ersten Batches werden als bereits ausgeführte Referenz
markiert; im Folge-Batch wurden PINHOLE und OPENCV als Original-GS-Ablationen
bei 2 und 5 FPS sowie drei Auflösungen getestet. Die SuGaR-Coarse-Auswertung
steht als separate, inzwischen abgeschlossene Folgematrix
\texttt{matrix\_sugar\_followup\_12} daneben.

\begin{table}[H]
\centering
\small
\begin{tabular}{lrrr}
\toprule
Batch & Versuche & vollständig erfolgreich & unvollständig\\
\midrule
matrix\_full\_pipe & 6 & 4 & 2\\
matrix\_rest & 18 & 12 & 6\\
Gesamt & 24 & 16 & 8\\
\bottomrule
\end{tabular}
\caption{Status der für die PA zusammengeführten Matrixbatches.}
\end{table}

\begin{figure}[H]
\centering
\includegraphics[width=\textwidth]{matrix_status_table.pdf}
\caption{Statusmatrix der archivierten Versuche.}
\end{figure}

Die sechs historischen SuGaR-Coarse-Läufe des Batches
\texttt{matrix\_rest} besitzen teilweise eine STS-Baseline-Metrik, sind aber
wegen des damaligen technischen Importfehlers im SuGaR-Render-Helfer nicht als
erfolgreiche SuGaR-Ergebnisse zu werten. Die Statusauswertung trennt daher
technische Vollständigkeit und vorhandene Zwischenmetriken. Die vollständige
STS-Übersicht liegt in \path{docs/grafiken/matrix_overview_table.pdf};
die kombinierte Overview der erfolgreichen STS-/A-Routen und SuGaR-Coarse-
Läufe liegt in \path{docs/grafiken/matrix_combined_overview_table.pdf};
die neue SuGaR-Coarse-Übersicht liegt in
\path{docs/grafiken/matrix_sugar_overview_table.pdf}. Die gemeinsame
CSV-Quelle liegt in \path{docs/grafiken/matrix_thesis_data.csv}; die
Spalte \texttt{model\_stage} trennt STS und SuGaR-Coarse.

\subsection*{Abgeschlossene SuGaR-Coarse-Folgematrix}

Die Konfiguration \path{tools/experiment_matrix_sugar.tsv} plant zwölf
neue Läufe: \path{SIMPLE_RADIAL/sugar_coarse} und
\path{OPENCV/sugar_coarse} bei drei Auflösungen und zwei FPS-Profilen. Der
Render-Helfer wurde vor dieser Prüfung durch einen expliziten
\texttt{/opt/sugar}-Importpfad und eine HWC-zu-CHW-Konvertierung korrigiert.
Ein archivierter 5-FPS-720p-Checkpoint ließ sich damit bereits auf 30
Testansichten rendern und objektmaskiert auswerten.

Die zwölf Läufe unter \path{data/10_runs/matrix_sugar_followup_12/}
sind abgeschlossen. Alle zwölf Manifeste tragen den Status \texttt{success}.
Die Coarse-Meshgewinnung, das SuGaR-Coarse-Rendering, die objektmaskierte
Metrikberechnung und das Postprocessing wurden für beide Kameramodelle sowie
alle Auflösungs-/FPS-Kombinationen erfolgreich archiviert. Die
\texttt{sugar\_refined\_masked.json}-Dateien weisen dagegen ausdrücklich
\texttt{no refined.ply exported} aus; die Folgematrix ist daher ein
SuGaR-Coarse-Nachweis und kein vollständiger SuGaR-Refinement-Vergleich.

\begin{table}[H]
\centering
\small
\begin{tabularx}{\textwidth}{l l rrr}
\toprule
Kameramodell & Auflösung & PSNR$\uparrow$ & SSIM$\uparrow$ & LPIPS$\downarrow$\\
\midrule
SIMPLE\_RADIAL & 720p & 21,380 & 0,764 & 0,1534\\
SIMPLE\_RADIAL & QHD  & 21,687 & 0,761 & 0,1407\\
SIMPLE\_RADIAL & Low  & 21,887 & 0,752 & 0,1261\\
OPENCV         & 720p & 21,453 & 0,764 & 0,1531\\
OPENCV         & QHD  & 21,624 & 0,764 & 0,1367\\
OPENCV         & Low  & 21,808 & 0,749 & 0,1255\\
\bottomrule
\end{tabularx}
\caption{SuGaR-Coarse-Mittelwerte aus 2 und 5 FPS der abgeschlossenen
Folgematrix. Die Werte sind objektmaskierte Ansichtsmetriken und kein
Geometrienachweis.}
\label{tab:sugar-followup-average}
\end{table}

Die vollständigen Einzelwerte und die getrennten Darstellungen für Status,
PSNR, SSIM, LPIPS und Zeit/Qualität liegen in den digitalen Anlagen unter
\path{docs/grafiken/matrix_sugar_*}. Alle zwölf Läufe verwendeten den
festen Eval-Split; bei 2 FPS wurden zwölf und bei 5 FPS dreißig gültige
Evaluationsansichten ausgewertet. Die Laufzeit-/Qualitätsgrafik nutzt die
archivierte vollständige Pipelinezeit und ist wie bei der STS-Matrix nur ein
sekundärer Screening-Index.

\clearpage
\begin{figure}[H]
\centering
\includegraphics[width=\textwidth]{matrix_sugar_status_table.pdf}
\caption{Statusübersicht der zwölf abgeschlossenen SuGaR-Coarse-Folgeläufe.}
\end{figure}

\clearpage
\begin{figure}[H]
\centering
\includegraphics[width=\textwidth]{matrix_sugar_psnr_table.pdf}
\caption{Objektmaskiertes PSNR der SuGaR-Coarse-Folgematrix.}
\end{figure}

\clearpage
\begin{figure}[H]
\centering
\includegraphics[width=\textwidth]{matrix_sugar_ssim_table.pdf}
\caption{Objektmaskiertes SSIM der SuGaR-Coarse-Folgematrix.}
\end{figure}

\clearpage
\begin{figure}[H]
\centering
\includegraphics[width=\textwidth]{matrix_sugar_lpips_table.pdf}
\caption{Objektmaskiertes LPIPS der SuGaR-Coarse-Folgematrix.}
\end{figure}

\clearpage
\begin{figure}[H]
\centering
\includegraphics[width=\textwidth]{matrix_sugar_time_quality_table.pdf}
\caption{Sekundäres Zeit-/Qualitätsverhältnis der SuGaR-Coarse-Folgematrix.}
\end{figure}

\clearpage
\subsection*{Regeln für die Auswertung}

\begin{itemize}
  \item \texttt{sts\_masked.json} ist die STS-Baseline vor dem separaten
	  SuGaR-Rendering und kein SuGaR-Ergebnis.
  \item Ein Average wird nur gebildet, wenn 2-FPS- und 5-FPS-Lauf derselben
	  Konfiguration vollständig erfolgreich sind.
  \item Fehlgeschlagene Läufe bleiben im Statusbericht und erhalten keinen
	  künstlich schlechten Metrikwert.
  \item Bildmetriken werden nur innerhalb nichtleerer Objektmasken berechnet.
  \item Native Auflösungen sind wegen Downsampling nicht unmittelbar als
	  Geometriequalität vergleichbar.
\end{itemize}
